% Options for packages loaded elsewhere
\PassOptionsToPackage{unicode}{hyperref}
\PassOptionsToPackage{hyphens}{url}
%
\documentclass[
]{article}
\usepackage{lmodern}
\usepackage{amssymb,amsmath}
\usepackage{ifxetex,ifluatex}
\ifnum 0\ifxetex 1\fi\ifluatex 1\fi=0 % if pdftex
  \usepackage[T1]{fontenc}
  \usepackage[utf8]{inputenc}
  \usepackage{textcomp} % provide euro and other symbols
\else % if luatex or xetex
  \usepackage{unicode-math}
  \defaultfontfeatures{Scale=MatchLowercase}
  \defaultfontfeatures[\rmfamily]{Ligatures=TeX,Scale=1}
\fi
% Use upquote if available, for straight quotes in verbatim environments
\IfFileExists{upquote.sty}{\usepackage{upquote}}{}
\IfFileExists{microtype.sty}{% use microtype if available
  \usepackage[]{microtype}
  \UseMicrotypeSet[protrusion]{basicmath} % disable protrusion for tt fonts
}{}
\makeatletter
\@ifundefined{KOMAClassName}{% if non-KOMA class
  \IfFileExists{parskip.sty}{%
    \usepackage{parskip}
  }{% else
    \setlength{\parindent}{0pt}
    \setlength{\parskip}{6pt plus 2pt minus 1pt}}
}{% if KOMA class
  \KOMAoptions{parskip=half}}
\makeatother
\usepackage{xcolor}
\IfFileExists{xurl.sty}{\usepackage{xurl}}{} % add URL line breaks if available
\IfFileExists{bookmark.sty}{\usepackage{bookmark}}{\usepackage{hyperref}}
\hypersetup{
  pdftitle={Data Handling: Import, Cleaning and Visualisation},
  pdfauthor={Prof.~Dr.~Ulrich Matter},
  hidelinks,
  pdfcreator={LaTeX via pandoc}}
\urlstyle{same} % disable monospaced font for URLs
\usepackage[margin=1in]{geometry}
\usepackage{color}
\usepackage{fancyvrb}
\newcommand{\VerbBar}{|}
\newcommand{\VERB}{\Verb[commandchars=\\\{\}]}
\DefineVerbatimEnvironment{Highlighting}{Verbatim}{commandchars=\\\{\}}
% Add ',fontsize=\small' for more characters per line
\usepackage{framed}
\definecolor{shadecolor}{RGB}{248,248,248}
\newenvironment{Shaded}{\begin{snugshade}}{\end{snugshade}}
\newcommand{\AlertTok}[1]{\textcolor[rgb]{0.94,0.16,0.16}{#1}}
\newcommand{\AnnotationTok}[1]{\textcolor[rgb]{0.56,0.35,0.01}{\textbf{\textit{#1}}}}
\newcommand{\AttributeTok}[1]{\textcolor[rgb]{0.77,0.63,0.00}{#1}}
\newcommand{\BaseNTok}[1]{\textcolor[rgb]{0.00,0.00,0.81}{#1}}
\newcommand{\BuiltInTok}[1]{#1}
\newcommand{\CharTok}[1]{\textcolor[rgb]{0.31,0.60,0.02}{#1}}
\newcommand{\CommentTok}[1]{\textcolor[rgb]{0.56,0.35,0.01}{\textit{#1}}}
\newcommand{\CommentVarTok}[1]{\textcolor[rgb]{0.56,0.35,0.01}{\textbf{\textit{#1}}}}
\newcommand{\ConstantTok}[1]{\textcolor[rgb]{0.00,0.00,0.00}{#1}}
\newcommand{\ControlFlowTok}[1]{\textcolor[rgb]{0.13,0.29,0.53}{\textbf{#1}}}
\newcommand{\DataTypeTok}[1]{\textcolor[rgb]{0.13,0.29,0.53}{#1}}
\newcommand{\DecValTok}[1]{\textcolor[rgb]{0.00,0.00,0.81}{#1}}
\newcommand{\DocumentationTok}[1]{\textcolor[rgb]{0.56,0.35,0.01}{\textbf{\textit{#1}}}}
\newcommand{\ErrorTok}[1]{\textcolor[rgb]{0.64,0.00,0.00}{\textbf{#1}}}
\newcommand{\ExtensionTok}[1]{#1}
\newcommand{\FloatTok}[1]{\textcolor[rgb]{0.00,0.00,0.81}{#1}}
\newcommand{\FunctionTok}[1]{\textcolor[rgb]{0.00,0.00,0.00}{#1}}
\newcommand{\ImportTok}[1]{#1}
\newcommand{\InformationTok}[1]{\textcolor[rgb]{0.56,0.35,0.01}{\textbf{\textit{#1}}}}
\newcommand{\KeywordTok}[1]{\textcolor[rgb]{0.13,0.29,0.53}{\textbf{#1}}}
\newcommand{\NormalTok}[1]{#1}
\newcommand{\OperatorTok}[1]{\textcolor[rgb]{0.81,0.36,0.00}{\textbf{#1}}}
\newcommand{\OtherTok}[1]{\textcolor[rgb]{0.56,0.35,0.01}{#1}}
\newcommand{\PreprocessorTok}[1]{\textcolor[rgb]{0.56,0.35,0.01}{\textit{#1}}}
\newcommand{\RegionMarkerTok}[1]{#1}
\newcommand{\SpecialCharTok}[1]{\textcolor[rgb]{0.00,0.00,0.00}{#1}}
\newcommand{\SpecialStringTok}[1]{\textcolor[rgb]{0.31,0.60,0.02}{#1}}
\newcommand{\StringTok}[1]{\textcolor[rgb]{0.31,0.60,0.02}{#1}}
\newcommand{\VariableTok}[1]{\textcolor[rgb]{0.00,0.00,0.00}{#1}}
\newcommand{\VerbatimStringTok}[1]{\textcolor[rgb]{0.31,0.60,0.02}{#1}}
\newcommand{\WarningTok}[1]{\textcolor[rgb]{0.56,0.35,0.01}{\textbf{\textit{#1}}}}
\usepackage{graphicx}
\makeatletter
\def\maxwidth{\ifdim\Gin@nat@width>\linewidth\linewidth\else\Gin@nat@width\fi}
\def\maxheight{\ifdim\Gin@nat@height>\textheight\textheight\else\Gin@nat@height\fi}
\makeatother
% Scale images if necessary, so that they will not overflow the page
% margins by default, and it is still possible to overwrite the defaults
% using explicit options in \includegraphics[width, height, ...]{}
\setkeys{Gin}{width=\maxwidth,height=\maxheight,keepaspectratio}
% Set default figure placement to htbp
\makeatletter
\def\fps@figure{htbp}
\makeatother
\setlength{\emergencystretch}{3em} % prevent overfull lines
\providecommand{\tightlist}{%
  \setlength{\itemsep}{0pt}\setlength{\parskip}{0pt}}
\setcounter{secnumdepth}{-\maxdimen} % remove section numbering
\usepackage[T1]{fontenc}
\usepackage{hyperref}

\title{Data Handling: Import, Cleaning and Visualisation}
\usepackage{etoolbox}
\makeatletter
\providecommand{\subtitle}[1]{% add subtitle to \maketitle
  \apptocmd{\@title}{\par {\large #1 \par}}{}{}
}
\makeatother
\subtitle{Exercise 5: Applied Data Analysis with R}
\author{Prof.~Dr.~Ulrich Matter}
\date{09/12/2021}

\begin{document}
\maketitle

\hypertarget{exercises}{%
\section[Exercises]{\texorpdfstring{Exercises\footnote{Adopted from
  \url{https://r4ds.had.co.nz/tidy-data.html}.}}{Exercises}}\label{exercises}}

\hypertarget{exercise-a}{%
\subsection{Exercise A}\label{exercise-a}}

Load \texttt{tidyverse} and read \texttt{stocks.csv} into RStudio (store
it in variable \texttt{stocks}). The dataset is currently in long
format. Use \texttt{pivot\_wider()} to transform it into wide format
(separate columns for different years). Store the resulting
\texttt{tibble}/\texttt{data.frame} in a variable called
\texttt{stocks2}. Then, run the following code.

\begin{Shaded}
\begin{Highlighting}[]
\KeywordTok{pivot\_longer}\NormalTok{(stocks2,  }\StringTok{\textasciigrave{}}\DataTypeTok{2015}\StringTok{\textasciigrave{}}\OperatorTok{:}\StringTok{\textasciigrave{}}\DataTypeTok{2016}\StringTok{\textasciigrave{}}\NormalTok{, }\DataTypeTok{names\_to =} \StringTok{"year"}\NormalTok{, }\DataTypeTok{values\_to =} \StringTok{"return"}\NormalTok{) }\OperatorTok{==}\StringTok{ }\NormalTok{stocks}
\end{Highlighting}
\end{Shaded}

Why is the result not identical to the original dataset
(\texttt{stocks})? In other words, why are \texttt{pivot\_longer()} and
\texttt{pivot\_wider()} not perfectly symmetrical? (Hint: look at the
variable types and think about column names.)

\begin{Shaded}
\begin{Highlighting}[]
\CommentTok{\# steps as instructed}
\KeywordTok{library}\NormalTok{(tidyverse)}
\NormalTok{stocks \textless{}{-}}\StringTok{ }\KeywordTok{read\_csv}\NormalTok{(}\StringTok{"../data/stocks.csv"}\NormalTok{)}
\NormalTok{stocks2 \textless{}{-}}\StringTok{ }\KeywordTok{pivot\_wider}\NormalTok{(stocks, }\DataTypeTok{names\_from =}\NormalTok{ year, }\DataTypeTok{values\_from =}\NormalTok{ return)}

\NormalTok{stocks\_wide\_long \textless{}{-}}\StringTok{ }\KeywordTok{pivot\_longer}\NormalTok{(stocks2, }\StringTok{\textasciigrave{}}\DataTypeTok{2015}\StringTok{\textasciigrave{}}\OperatorTok{:}\StringTok{\textasciigrave{}}\DataTypeTok{2016}\StringTok{\textasciigrave{}}\NormalTok{, }\DataTypeTok{names\_to =} \StringTok{"year"}\NormalTok{, }
                                 \DataTypeTok{values\_to =} \StringTok{"return"}\NormalTok{) }

\CommentTok{\# test if symmetrical}
\NormalTok{stocks\_wide\_long }\OperatorTok{==}\StringTok{ }\NormalTok{stocks}
\end{Highlighting}
\end{Shaded}

\begin{verbatim}
##       half  year return
## [1,] FALSE FALSE   TRUE
## [2,] FALSE FALSE  FALSE
## [3,] FALSE FALSE  FALSE
## [4,] FALSE FALSE   TRUE
\end{verbatim}

\begin{Shaded}
\begin{Highlighting}[]
\CommentTok{\# another way to test whether the datasets are exactly equal}
\KeywordTok{identical}\NormalTok{(stocks\_wide\_long, stocks)}
\end{Highlighting}
\end{Shaded}

\begin{verbatim}
## [1] FALSE
\end{verbatim}

In order to evaluate the differences between the two datasets
\texttt{stocks} and \texttt{stocks\_wide\_long}, take into account the
variable types and the ordering of the columns.

\begin{Shaded}
\begin{Highlighting}[]
\CommentTok{\# investigate the variable types}
\KeywordTok{str}\NormalTok{(stocks\_wide\_long)}
\KeywordTok{str}\NormalTok{(stocks)}
\end{Highlighting}
\end{Shaded}

The variables are not of the same types. They are also given in
different order.

\hypertarget{exercise-b}{%
\subsection{Exercise B}\label{exercise-b}}

Load \texttt{tidyverse} and read \texttt{countries.csv} into R (store it
in variable \texttt{countries}). Then, run the following code.

\begin{Shaded}
\begin{Highlighting}[]
\KeywordTok{pivot\_longer}\NormalTok{(countries, }\DecValTok{1999}\NormalTok{, }\DecValTok{2000}\NormalTok{, }\DataTypeTok{names\_to =} \StringTok{"year"}\NormalTok{, }\DataTypeTok{values\_to =} \StringTok{"cases"}\NormalTok{)}
\end{Highlighting}
\end{Shaded}

Why does this code fail?

\begin{Shaded}
\begin{Highlighting}[]
\CommentTok{\# read the data}
\NormalTok{countries \textless{}{-}}\StringTok{ }\KeywordTok{read\_csv}\NormalTok{(}\StringTok{"../data/countries.csv"}\NormalTok{)}

\CommentTok{\# try the example code}
\KeywordTok{pivot\_longer}\NormalTok{(countries, }\KeywordTok{c}\NormalTok{(}\DecValTok{1999}\NormalTok{, }\DecValTok{2000}\NormalTok{), }\DataTypeTok{names\_to =}  \StringTok{"year"}\NormalTok{, }\DataTypeTok{values\_to =} \StringTok{"cases"}\NormalTok{)}
\end{Highlighting}
\end{Shaded}

The code fails because numeric values cannot be used as variable names.

\begin{Shaded}
\begin{Highlighting}[]
\CommentTok{\# fix the problem }
\KeywordTok{pivot\_longer}\NormalTok{(countries, }\KeywordTok{c}\NormalTok{(}\StringTok{\textasciigrave{}}\DataTypeTok{1999}\StringTok{\textasciigrave{}}\NormalTok{, }\StringTok{\textasciigrave{}}\DataTypeTok{2000}\StringTok{\textasciigrave{}}\NormalTok{), }\DataTypeTok{names\_to =} \StringTok{"year"}\NormalTok{, }\DataTypeTok{values\_to =} \StringTok{"cases"}\NormalTok{)}
\end{Highlighting}
\end{Shaded}

\hypertarget{exercise-c}{%
\subsection{Exercise C}\label{exercise-c}}

Load \texttt{tidyverse} and read \texttt{people.csv} into R (store it in
variable \texttt{people}). Try to spread \texttt{people}. Why does the
spreading fail? How could you add a new column to fix the problem?

\begin{Shaded}
\begin{Highlighting}[]
\CommentTok{\# read the data}
\NormalTok{people \textless{}{-}}\StringTok{ }\KeywordTok{read\_csv}\NormalTok{(}\StringTok{"../data/people.csv"}\NormalTok{)}
\KeywordTok{pivot\_wider}\NormalTok{(people, }\DataTypeTok{names\_from =}\NormalTok{ key, }\DataTypeTok{values\_from =}\NormalTok{  value)}

\CommentTok{\# look at the dataset}
\KeywordTok{view}\NormalTok{(people)}
\KeywordTok{str}\NormalTok{(people)}
\end{Highlighting}
\end{Shaded}

The issue is that the key-variable -- the identifier -- contains
duplicate values.

\begin{Shaded}
\begin{Highlighting}[]
\CommentTok{\# fix problem:}
\CommentTok{\# assume there are two Phillip Woods and mark them with a different identifier}
\NormalTok{people}\OperatorTok{$}\NormalTok{id \textless{}{-}}\StringTok{ }\KeywordTok{c}\NormalTok{(}\DecValTok{1}\NormalTok{,}\DecValTok{1}\NormalTok{,}\DecValTok{2}\NormalTok{,}\DecValTok{3}\NormalTok{,}\DecValTok{3}\NormalTok{)}
\KeywordTok{pivot\_wider}\NormalTok{(people, }\DataTypeTok{names\_from =}\NormalTok{ key, }\DataTypeTok{values\_from =}\NormalTok{  value)}
\end{Highlighting}
\end{Shaded}

\hypertarget{exercise-d}{%
\subsection{Exercise D}\label{exercise-d}}

Read \texttt{preg.csv} into RStudio (store it in variable
\texttt{preg}). Tidy \texttt{preg}. Do you need to use
\texttt{pivot\_wider} or \texttt{pivot\_longer}? What are the variables?

\begin{Shaded}
\begin{Highlighting}[]
\CommentTok{\# read the data}
\NormalTok{preg \textless{}{-}}\StringTok{ }\KeywordTok{read\_csv}\NormalTok{(}\StringTok{"../data/preg.csv"}\NormalTok{)}

\CommentTok{\# look at the dataset}
\KeywordTok{view}\NormalTok{(preg)}

\CommentTok{\# use common sense (\textquotesingle{}heuristics\textquotesingle{}) to figure out what information are contained}
\KeywordTok{pivot\_longer}\NormalTok{(preg, }\KeywordTok{c}\NormalTok{(}\StringTok{"male"}\NormalTok{, }\StringTok{"female"}\NormalTok{), }\DataTypeTok{names\_to =} \StringTok{"gender"}\NormalTok{, }\DataTypeTok{values\_to =} \StringTok{"count"}\NormalTok{)}
\end{Highlighting}
\end{Shaded}

\hypertarget{additional-exercises-advanced}{%
\section{Additional Exercises
(advanced)}\label{additional-exercises-advanced}}

\hypertarget{exercise-f}{%
\subsection{Exercise F}\label{exercise-f}}

You want to investigate the relationship between GDP per capita and CO2
emissions per capita across countries. On Canvas, download the dataset
\texttt{co2gdp.csv}. This dataset has been downloaded from the website
``Our World in Data''
(\url{https://ourworldindata.org/grapher/co2-emissions-vs-gdp}).
Investigate the relationship between the dependent variable
\texttt{CO2\_percapita} and the independent variable
\texttt{GDP\_percapita} using a linear regression of the following form:

\[\text{CO2_percapita}_i = \alpha + \beta \text{GDP_percapita}_i + u_i.\]

Estimate this linear regression by creating your own function
\texttt{myols} that takes the following three arguments: (\(x\)) for the
dependent variable, (\(y\)) for the independent variable, and
\textit{matrix}, which takes \texttt{TRUE} if the function estimates the
coefficients using matrix algebra, and \texttt{FALSE} if the function
estimates the coefficients using the sum operators. Compare the
coefficients with the \texttt{R} built-in function \texttt{lm()}.
Finally, comment and interpret the coefficients.

\textbf{Hint 1} (linear regression with sum operator): the slope
parameter of the linear regression is expressed as the following
formula:
\[\hat{\beta} = \frac{\sum{x_{i}y_{i}}-N\bar{y}\bar{x}}{\sum{x_i^2}-N\bar{x}^2},\]
and the constant is given by this formula:
\[\hat{\alpha}=\bar{y}-\hat{\beta}\bar{x},\] where \(\bar{x}\) gives the
mean of \(x\) and \(\bar{y}\) gives the mean of \(y\).

\textbf{Hint 2} (linear regression with matrix algebra): the formula for
the linear regression in matrix algebra is
\[\hat{\bf{\beta}} = (\bf{X'X})^{-1} \bf{Xy},\] where you find the
inverse of a matrix with the command `solve()', the cross-product of two
matrices with \texttt{crossprod(X,Y)} \(= t(X)\%^{*}\% Y\). Also, note
that you need a column of 1s in your matrix of covariates to compute the
intercept.

(\textbf{Note}: this exercise is about programming and implementing the
formula of the linear regression. You don't need to understand the
mechanics of linear regression for the exam. This will be covered next
semester in the course \texttt{Data\ Analytics\ II}.)

\begin{Shaded}
\begin{Highlighting}[]
\CommentTok{\# Create a function for linear regression}
\NormalTok{myols \textless{}{-}}\StringTok{ }\ControlFlowTok{function}\NormalTok{(y, x, }\DataTypeTok{matrix =}\NormalTok{ F) \{ }\CommentTok{\# The function is called "myols". It takes three arguments: y, x, and TRUE or FALSE for matrix.}
  
  \ControlFlowTok{if}\NormalTok{ (}\KeywordTok{isFALSE}\NormalTok{(matrix))\{ }\CommentTok{\# if matrix is FALSE, we apply the formula for the slope coefficient}
  
  \CommentTok{\# Apply formula for the slope coefficient}
\NormalTok{  N \textless{}{-}}\StringTok{ }\KeywordTok{length}\NormalTok{(y)}
    
\NormalTok{  beta \textless{}{-}}\StringTok{ }\NormalTok{(}\KeywordTok{sum}\NormalTok{(y}\OperatorTok{*}\NormalTok{x) }\OperatorTok{{-}}\StringTok{ }\NormalTok{N}\OperatorTok{*}\KeywordTok{mean}\NormalTok{(x)}\OperatorTok{*}\KeywordTok{mean}\NormalTok{(y)) }\OperatorTok{/}\StringTok{ }\NormalTok{(}\KeywordTok{sum}\NormalTok{(x}\OperatorTok{\^{}}\DecValTok{2}\NormalTok{)}\OperatorTok{{-}}\NormalTok{N}\OperatorTok{*}\KeywordTok{mean}\NormalTok{(x)}\OperatorTok{\^{}}\DecValTok{2}\NormalTok{) }\CommentTok{\# apply the formula given in the exercise}
  
  \CommentTok{\# Apply formula for the constant}
\NormalTok{  alpha \textless{}{-}}\StringTok{ }\KeywordTok{mean}\NormalTok{(y)}\OperatorTok{{-}}\NormalTok{beta}\OperatorTok{*}\KeywordTok{mean}\NormalTok{(x)}
  
  \KeywordTok{return}\NormalTok{(}\KeywordTok{list}\NormalTok{(}\StringTok{"alpha"}\NormalTok{ =}\StringTok{ }\NormalTok{alpha, }\StringTok{"beta"}\NormalTok{ =}\StringTok{ }\NormalTok{beta)) }\CommentTok{\# The function returns a list with two named elements}
\NormalTok{  \}}

  \ControlFlowTok{if}\NormalTok{ (}\KeywordTok{isTRUE}\NormalTok{(matrix))\{ }\CommentTok{\# if matrix is TRUE, we apply the formula for the slope coefficient in matrix notation}
    
    \CommentTok{\# Compute cross products and inverse}
\NormalTok{    x \textless{}{-}}\StringTok{ }\KeywordTok{cbind}\NormalTok{(}\DecValTok{1}\NormalTok{, x) }\CommentTok{\# we first need to add a column of 1 for the constant}
\NormalTok{    XXi \textless{}{-}}\StringTok{ }\KeywordTok{solve}\NormalTok{(}\KeywordTok{crossprod}\NormalTok{(x,x)) }\CommentTok{\# we apply the formula given in the exercise}
\NormalTok{    Xy \textless{}{-}}\StringTok{ }\KeywordTok{crossprod}\NormalTok{(x, y) }\CommentTok{\# idem}
          
    \KeywordTok{return}\NormalTok{(XXi }\OperatorTok{\%*\%}\StringTok{ }\NormalTok{Xy) }\CommentTok{\# idem}
\NormalTok{  \}}
\NormalTok{\} }\CommentTok{\# The function is closed.}


\CommentTok{\# Import data}
\NormalTok{co2 \textless{}{-}}\StringTok{ }\KeywordTok{read\_csv}\NormalTok{(}\StringTok{"../data/co2gdp.csv"}\NormalTok{)}

\CommentTok{\# Run the function}
\KeywordTok{myols}\NormalTok{(}\DataTypeTok{y =}\NormalTok{ co2}\OperatorTok{$}\NormalTok{CO2\_percapita, }\DataTypeTok{x =}\NormalTok{ co2}\OperatorTok{$}\NormalTok{GDP\_percapita, }\DataTypeTok{matrix =}\NormalTok{ F)}
\end{Highlighting}
\end{Shaded}

\begin{verbatim}
## $alpha
## [1] 0.4931492
## 
## $beta
## [1] 0.000379723
\end{verbatim}

\begin{Shaded}
\begin{Highlighting}[]
\KeywordTok{myols}\NormalTok{(}\DataTypeTok{y =}\NormalTok{ co2}\OperatorTok{$}\NormalTok{CO2\_percapita, }\DataTypeTok{x =}\NormalTok{ co2}\OperatorTok{$}\NormalTok{GDP\_percapita, }\DataTypeTok{matrix =}\NormalTok{ T)}
\end{Highlighting}
\end{Shaded}

\begin{verbatim}
##          [,1]
##   0.493149193
## x 0.000379723
\end{verbatim}

\begin{Shaded}
\begin{Highlighting}[]
\CommentTok{\# Check with lm()}
\KeywordTok{lm}\NormalTok{(CO2\_percapita }\OperatorTok{\textasciitilde{}}\StringTok{ }\NormalTok{GDP\_percapita, }\DataTypeTok{data =}\NormalTok{ co2)}
\end{Highlighting}
\end{Shaded}

\begin{verbatim}
## 
## Call:
## lm(formula = CO2_percapita ~ GDP_percapita, data = co2)
## 
## Coefficients:
##   (Intercept)  GDP_percapita  
##     0.4931492      0.0003797
\end{verbatim}

\begin{Shaded}
\begin{Highlighting}[]
\CommentTok{\# Graph}
\KeywordTok{ggplot}\NormalTok{(co2, }\KeywordTok{aes}\NormalTok{(}\DataTypeTok{x =}\NormalTok{ GDP\_percapita, }\DataTypeTok{y =}\NormalTok{ CO2\_percapita)) }\OperatorTok{+}
\StringTok{  }\KeywordTok{geom\_point}\NormalTok{() }\OperatorTok{+}
\StringTok{  }\KeywordTok{geom\_smooth}\NormalTok{(}\DataTypeTok{method =} \StringTok{"lm"}\NormalTok{, }\DataTypeTok{se =}\NormalTok{ F)}
\end{Highlighting}
\end{Shaded}

\includegraphics{05_data_analysis_files/figure-latex/unnamed-chunk-11-1.pdf}

\hypertarget{references}{%
\section{References}\label{references}}

\end{document}
